\Article{25}
Neben den angeführten Versuchen haben wir viele weitere ähnliche Versuche durchgeführt, die früheren allerdings mit erheblich geringerer Sorgfalt. Dennoch ist es nützlich, ihre Ergebnisse hier zusammenzustellen. Ausgenommen sind die Ergebnisse, die vor der Aufstellung der feineren Apparate mit einfacheren Hilfsmitteln an Nadeln sehr unterschiedlicher Abmessungen gewonnen wurden, obwohl \textit{alle} wenigstens eine Annäherung an den wahren Wert ergaben. Bei den wiederholten Versuchen erhielten wir nacheinander folgende Werte von $T$.
\SourceLeaf{119}
\begin{center}
\begin{tabular}{c|l|l}
Nummer&Datum, $1832$&$T$\\
I&Mai $21$&$1{,}7820$\\
II&Mai $24$&$1{,}7694$\\
III&Juni $4$&$1{,}7713$\\
IV&Juni $24$—$28$&$1{,}7625$\\
V&Juli $23$, $24$&$1{,}7826$\\
VI&Juli $25$, $26$&$1{,}7845$\\
VII&Sept. $9$&$1{,}7764$\\
VIII&Sept. $18$&$1{,}7821$\\
IX&Sept. $27$&$1{,}7965$\\
X&Okt. $15$&$1{,}7860$
\end{tabular}
\end{center}
Die Versuche V—IX wurden alle am selben Ort durchgeführt, die Versuche I—IV dagegen an anderen Orten. Versuch X ist eigentlich ein gemischter Fall, da die Ablenkungen am gewohnten Ort, die Schwingungen dagegen anderswo beobachtet wurden. Bei den Versuchen VII und VIII wurde fast dieselbe Sorgfalt aufgewandt, bei IV, V, VI, X etwas weniger und bei I—III erheblich weniger. Für die Versuche I—VIII wurden verschiedene Nadeln von nahezu gleichem Gewicht und gleicher Länge verwendet (das Gewicht lag zwischen $400$ und $440$ Gramm). Bei Versuch X kam dagegen eine Nadel mit einem Gewicht von $1062$ Gramm und einer Länge von $485$ Millimetern zum Einsatz. Versuch IX wurde nur unternommen, um zu sehen, welcher Grad an Genauigkeit mit einer eher kleinen Nadel erreichbar ist. Die verwendete Nadel wog nur $58$ Gramm, doch wurde im Übrigen nicht weniger Sorgfalt aufgewandt als bei den Versuchen VII und VIII. Zweifellos wird die Genauigkeit der Beobachtungen merklich zunehmen, wenn noch schwerere Nadeln verwendet werden, etwa solche mit einem Gewicht von bis zu $2000$ oder $3000$ Gramm.

